\section{Introduction}
\label{sec:intro}

Persistent homology (PH) is now a standard tool for characterizing the shape of high-dimensional point clouds, and it has been applied to bulk and single-cell omics data to argue for the presence of loops, branchings, and other topological features that a purely linear or clustering-based analysis would miss~\cite{placeholder-ph-omics}. In a typical workflow, samples are embedded in gene or feature space, reduced via principal component analysis (PCA) to a tractable number of dimensions, and a Vietoris--Rips filtration is built on the resulting point cloud; persistent first-homology ($H_1$) features --- one-dimensional ``loops'' --- are then reported as evidence of nontrivial biological structure: cyclic trajectories, multi-way transitions between cell states, or recurrent sub-populations that PCA has ``revealed.''

This inference has a well-known blind spot. PCA is not merely a projection; on any point cloud, in any number of dimensions, projecting to a lower-dimensional subspace changes the pairwise distance structure and can inflate or deflate persistence values in ways that have nothing to do with the data's substantive structure. When PCA-space $H_1$ persistence is reported as evidence of biological topology, the comparison is almost never made against a null model constructed from the \emph{same} preprocessing pipeline applied to unstructured data. Without that comparison, a claim of the form ``PCA reveals structure'' is unfalsifiable: an equally large signal could arise from resampled noise pushed through the identical highly-variable-gene selection, PCA, and Vietoris--Rips steps.

This project tests two precise, pre-registered, and mutually distinguishable hypotheses against that gap.

\begin{description}
\item[$H_0^{\text{stat}}$ (signal detection).] The maximum $H_1$ persistence observed in real omics data significantly exceeds that of two matched null models --- an i.i.d.\ Gaussian null matched in ambient dimension and sample count, and a pipeline-symmetric permutation null that shuffles each gene's values independently across samples before re-running the identical highly-variable-gene selection, PCA, and filtration steps --- in \emph{both} raw feature space and PCA-reduced space, with effect size $z>3.0$ against both nulls in each space.
\item[$H_1^{\text{stat}}$ (PCA-attribution).] The literal claim tested is narrower and more adversarial to the naive interpretation: does PCA inflate $\maxHone$ in real data by an amount that itself exceeds the inflation the same PCA step induces on null data? If PCA-driven inflation in noise is comparable to or larger than PCA-driven inflation in real signal, then citing a PCA-space $H_1$ feature as evidence of biological topology is not supported by the data --- the inflation itself carries no information.
\end{description}

These two hypotheses are logically independent: a dataset can pass $H_0^{\text{stat}}$ (real data's absolute $\maxHone$ exceeds null) while still failing to license a naive ``PCA reveals structure'' reading if $H_1^{\text{stat}}$'s attribution test shows the PCA step itself is not doing discriminative work. Distinguishing them is the central methodological contribution of this program.

The results reported here rest on three pieces of work, each pre-registered or protocol-locked before analysis: (1) a pilot study (GSE81089, NSCLC bulk RNA-seq) followed by three independent, pre-registered confirmatory replications spanning four tissue types and two additional omics modalities (DNA methylation, mass-spectrometry proteomics); (2) an 18-configuration ablation sweep over the pilot cohort's hyperparameters (highly-variable-gene count, PCA component count, permutation count, missing-value imputation rule) to establish the boundary of the pre-registered operating point's validity; and (3) a confound-attribution audit designed to test whether an $H_1$ signal that does survive both nulls is a genuine feature of the omics manifold or merely a re-encoding of the tumor/normal class-mean separation that a linear classifier can already recover at near-ceiling accuracy in every cohort tested.

We treat the third question --- confound attribution --- as open rather than resolved. The audit protocol has so far been executed to completion only on the pilot cohort, where it finds the topological signal is not attributable to the class-mean shift; whether this holds in the three replication cohorts is a live, unanswered question that this paper tests but does not yet close, and one of the confirmatory cohorts (GSE146889) already shows a structural reason for caution: its dominant $H_1$ loop is significantly enriched for tumor-status samples, unlike the pilot's loop, meaning confound-independence found in one cohort should not be assumed to transfer to another without direct testing. \TODO{final synthesis of confound-audit generalization across all four cohorts to be inserted once the remaining three audits complete}

The remainder of the paper proceeds as follows. \Cref{sec:methods} describes the persistent-homology and null-model pipeline, the five cohorts and their two additional omics modalities, the ablation sweep design, and the confound-attribution audit protocol. \Cref{sec:results} reports the cross-dataset $H_0^{\text{stat}}$/$H_1^{\text{stat}}$ replication results, the ablation sweep's hyperparameter-sensitivity boundary, and the pilot-cohort confound-audit finding. \Cref{sec:discussion} synthesizes these results into a scope statement for what this program does and does not establish. \TODO{Discussion section to be finalized once confound-audit generalization results are available}
